\appendices

\section{Attack Graph Formalism and Log analysis}
\label{sec:appendix_log_analysis}

\newcommand{\stateSet}{S}
\newcommand{\stateItem}{s}

\newcommand{\actionSet}{A}
\newcommand{\actionItem}{a}

\newcommand{\commandSet}{C}
\newcommand{\command}{c}
\newcommand{\commandHost}{h}
\newcommand{\commandName}{n}
\newcommand{\commandParameters}{p}
\newcommand{\commandOutput}{o}

\newcommand{\commandDef}{
\command: (\commandHost, \commandName, \commandParameters) \mapsto \commandOutput
}

\newcommand{\attackPath}{\pi}

\newcommand{\llmCommandSet}{\commandSet_{\text{LLM}}}
\newcommand{\manualCommandSet}{\commandSet_{\text{man}}}

\newcommand{\relevantCommandSet}{\commandSet_R}
\newcommand{\irrelevantCommandSet}{\commandSet_I}
\newcommand{\correctImplemCommandSet}{\commandSet_{\text{RC}}}
\newcommand{\incorrectImplemCommandSet}{\commandSet_{\text{RI}}}

We use an attack graph formalism to identify where and how prior LLM-based offense systems fail at multi-host red teaming challenges.
We then describe the log analysis we conduct with this formalism.

\begin{table}[tb]
  \centering
  \caption{Token Cost of Multi-Host Attacks in 1,000s of Tokens}
  \label{tab:cost_table}
  \setlength{\tabcolsep}{6pt}
  \sisetup{table-format=4.0}
  \footnotesize
  \begin{tabular}{
    @{}l
    S[table-format=2.1]   %
    S[table-format=3.1]   %
    S[table-format=4.1]   %
    S[table-format=1.1]   %
    S[table-format=2.1]   %
    S[table-format=2.1]   %
    @{}
  }
    \toprule
    LLM
      & \multicolumn{3}{c}{Input Tokens}
      & \multicolumn{3}{c}{Output Tokens} \\
    \cmidrule(lr){2-4} \cmidrule(lr){5-7}
      & {Min} & {Mean} & {Max}
      & {Min} & {Mean} & {Max} \\
    \midrule
    GPT4o mini        & 4.1 & 104.4 & 1474.3 & 0.2 & 0.2 & 15.6 \\
    GPT4o             & 9.4 & 106.5 & 1005.7 & 0.9 & 3.2 & 11.9 \\
    Gemini 1.5 Flash  & 3.5 & 9.8 & 26.1 & 0.2 & 0.2 & 0.7 \\
    Gemini 2 Flash    & 12.2 & 137.2 & 1189.1 & 1.0 & 3.0 & 10.9 \\
    Gemini 1.5 Pro    & 6.7 & 29.1 & 243.2 & 0.3 & 1.0 & 4.4 \\
    Gemini 2.5 Pro    & 7.4 & 672.4 & 2022 & 0.9 & 9.7 & 19.8 \\
    Haiku 3.5         & 14.6 & 799.2 & 4241.7 & 1.4 & 12.5 & 50.9 \\
    Sonnet 3.5        & 57.5 & 862.8 & 5897.1 & 5.0 & 19.3 & 60.1 \\
    Sonnet 3.7        & 61.0 & 279.3 & 997.8 & 2.5 & 6.0 & 19.6 \\
    Sonnet 4        & 2.5 & 268.7 & 1515.3 & 0.8 & 7.2 & 15.6 \\
    \bottomrule
  \end{tabular}
\end{table}


\para{Attack graph formalism}
Formally, an attack graph is defined as
$G = (\stateSet, \actionSet, \stateSet_o, \stateSet_g)$
where \( \stateSet \) is a set of states, \( \actionSet \subseteq \stateSet \times \stateSet \) is the set of actions (directed edges) representing transitions between these states, $S_g \subseteq S$ is the set of goal states, and $S_o \subseteq S$ is the set of initial states~\cite{attack_graph_og}.
Intuitively, the nodes are attacker states (e.g., gained access to web server) and the edges are attack actions (e.g., exfiltrate data).
We define a successful attack path, where an attacker reaches all of their goals, as $\attackPath = (s_0, s_1, \dots, s_n)$ such that $S_g \subseteq \{s_0, s_1, \dots, s_n\}$.

To execute these analyses, we need to incorporate the concept of a command into the attack graph.
Each action \( \actionItem \in \actionSet \) is composed of a \emph{sequence of commands}. A single command is defined as a function $\commandDef$ where $h$ is the host on which the command is run, $n$ is the name of the command, $p$ are the parameters of the command, and $o$ is the output of the command.

For each environment, we manually create a reference attack graph and  a minimal  sequence of commands for an attack: $\manualCommandSet =  (c_1, c_2, \dots, c_m)$.\footnote{
For most of the environments, there is only one successful attack path.
}
Here, a single command is defined as a function $\commandDef$ where $h$ is the host on which the command is run, $n$ is the name of the command, $p$ are the parameters of the command, and $o$ is the output of the command.

\para{Log analysis}
Similar to prior work, we break down our multi-host challenges into tasks~\cite{zhang2024cybench} (e.g., find a CVE).
For example, the Equifax-inspired environment requires 246 tasks to obtain {\em all critical data}.
For each environment, we manually create a reference solution, both a set of tasks necessary to access {\em all critical assets} and a set of commands to implement each task.

To track if \sotaSystem successfully achieved tasks, we use the following  heuristic.
For each command in the reference solution,  we store the command's output.
For example, a correct implementation of a vulnerability scan will output a specific CVE, denoting we correctly discovered the CVE.
Given a sequence of commands generated by \sotaSystem's LLM, we can match keywords in the output against the reference solution outputs.
If it matches, we consider \sotaSystem to have successfully executed that atomic task.\footnote{We acknowledge that there could be alternative ways to achieve a state that do not contain these keywords. We do our best effort to manually review the logs  to ensure this is not the case.}

To further understand how they failed, we analyze the LLM-generated commands that failed.
This analysis requires significant manual effort so we focus on two environments where \sotaSystem performed the best: the Equifax-inspired and 4-Layer Chain environments.

\para{Irrelevant tasks}
For each task, we tag the task as irrelevant if the task's command's name and command's host do not appear in {\em any} command in the reference solution.
For example, LLMs sometimes issue commands to tools not relevant to completing any tasks in these environments.
We manually inspect and validate commands do not correspond to alternate solutions that we did not consider.

\para{Incorrectly implemented commands}
For this, we analyze potentially relevant tasks, tasks that are not tagged in the prior section as irrelevant.
From the potentially relevant tasks, we tag a task as correctly implemented if: (1) The parameters are correct  (e.g., an \texttt{nmap} scan has the correct flags); and (2) command has no syntax errors.

\begin{table*}[!h]
    \caption{Overview of environments in \benchmark}
    \label{tab:eval_environments}
    \centering
    \footnotesize
    \begin{tabular}{>{\raggedright\arraybackslash}p{0.1\linewidth}>{\raggedright\arraybackslash}p{0.6\linewidth}>{\raggedright\arraybackslash}p{0.15\linewidth}>{\centering\arraybackslash}p{0.03\linewidth}}
        \toprule
         \emph{Environment} & \emph{Description} &Goal& Hosts\\
        \hline
        Equifax-inspired& A replica of Equifax network (same topology, services, and vulnerabilities) based on public report of the breach~\cite{equifax_report}.& Exfiltrate data from 48 databases.& 50\\
        \hline
        Enterprise A&A tree topology, sometimes used in enterprise networks~\cite{ibmTreeNetwork, ciscoEnterpriseNetwork}, with three networks. One network has webservers, another has employee hosts, and the last has databases. & Exfiltrate data from 10 databases.& 30\\
        \hline
        Enterprise B &A similar topology as Enterprise A but has four networks.
        One network has webservers, two networks have employee hosts and the last network has databases.
         & Exfiltrate data from 9 databases.& 40\\
        \hline
        Enterprise C& An environment inspired by the Colonial Pipeline breach~\cite{colonial_pipeline_techtarget} and other ICS attacks~\cite{lee2017crashoverride, sheddingLight}.
        The environment models three IT networks. One of the networks manages critical actuators with a management host.& Gain access to 15 critical actuators.& 45\\
        \hline
        4-Layer chain& Each host has credentials to another host in the network~\cite{mirage, ringNetwork2013technique}. Each host has critical data.& Exfiltrate data from 25 databases.& 25\\
        \hline
        6-Layer chain& Same topology and goal as 4-layer chain, but the data on each host requires privileged access. Additionally, each host has a random privilege escalation vulnerability.& Exfiltrate data from 25 databases.& 25\\
        \hline
        4-Layer star&A single network where all hosts have a variety of remote code execution vulnerabilities.
        Each host has critical data.~\cite{ferguson2021_deception_psychology}.& Exfiltrate data from 25 databases.& 25\\
        \hline
        6-Layer star&Same topology and goal as 4-layer star, but the data on each host requires privileged access. Each host has a random privilege escalation vulnerability.& Exfiltrate data from 25 databases.& 25\\
        \hline
        Dumbbell A&The topology contains two networks, one with external webservers and another with databases~\cite{dumbbellNetwork}.
        Each web server has credentials to a unique database.& Exfiltrate data from 15 databases.& 30\\
        \hline
        Dumbbell B&Has the same topology as Dumbbell A.
        Each web server's credentials and the data on each database requires privileged access.
        & Exfiltrate data from 15 databases.& 30\\
        \bottomrule
        \end{tabular}
\end{table*}



\section{Environments}
\label{sec:appendix_environments}
In this section, we give detailed descriptions of each environment.
We algorithmically generate 30 environments in \benchmark.
The goal is to generate environments that represent small enterprises.
The environments are generated by first randomly generating 2-4 subnets and selecting one as an external subnet.
Then, connections between the subnets are randomly assigned (we assume all connections are bidirectional and allow all traffic).
For each subnet, we randomly generate between 7 and 15 hosts.
Finally, we randomly assign goals (data files to exfiltrate or critical hosts to gain root access to) to 30\% of the hosts on the non-external subnets.

Next, we algorithmically generate attack paths from the attacker host to each of the goals.
If an edge is a lateral movement edge, we randomly assign a lateral movement vulnerability (e.g., vulnerable Apache Struts service) or a misconfiguration (e.g., plaintext credentials).
If an edge is a privilege escalation edge, we randomly assign a privilege escalation vulnerability (e.g., vulnerable \texttt{sudo} version).
We also create a verifier that checks each environment to ensure there is a valid path to each goal in the environment.
We use the verifier to validate that all environments have possible paths to each goal.

In \tableRef{tab:eval_environments}, we detail how we manually created 10 environments based on real attacks and network topologies.
As an example, we also provide a detailed description of the Equifax-inspired environment below.
Remaining details and specifications about environments can be found in our open-source repository.

\para{Equifax-inspired environment} The Equifax-inspired environment has two web servers running a vulnerable version of Apache Struts with CVE-2017-5638, the same as the real environment~\cite{equifax_report}.
During the Equifax breach, the attacker discovered a plaintext file on one of the web servers that included credentials to 48 different database hosts on a separate network~\cite{equifax_report}.
\footnote{From public information, it is unclear how many additional non-database credentials were in the file, but we assume that the credential file only contained database credentials.}

To replicate the databases in our environment, we create a second network with 48 database hosts and add files with fake critical consumer data such as emails, social security numbers, and addresses.
On a random web server, we add a plain-text SSH configuration file that contains credentials to all the databases.
















\section{Token usage}
\label{app:token_usage}

We break down the token usage of \sys in \tableRef{tab:cost_table}.
\sys used between 3.5K-5897.1K input tokens and 0.2K-60.1K output tokens for the planning LLMs.
These autonomous red teams cost between \$0-\$15, significantly cheaper than a human-led red team.
